\documentclass[
  bibliography=totoc,     % Literatur im Inhaltsverzeichnis
  captions=tableheading,  % Tabellenüberschriften
  titlepage=firstiscover, % Titelseite ist Deckblatt
]{scrartcl}

% Warnung, falls nochmal kompiliert werden muss
\usepackage[aux]{rerunfilecheck}

% unverzichtbare Mathe-Befehle
\usepackage{amsmath}
% viele Mathe-Symbole
\usepackage{amssymb}
% Erweiterungen für amsmath
\usepackage{mathtools}

% Fonteinstellungen
\usepackage{fontspec}
% Latin Modern Fonts werden automatisch geladen
% Alternativ zum Beispiel:
%\setromanfont{Libertinus Serif}
%\setsansfont{Libertinus Sans}
%\setmonofont{Libertinus Mono}

% Wenn man andere Schriftarten gesetzt hat,
% sollte man das Seiten-Layout neu berechnen lassen
\recalctypearea{}

% deutsche Spracheinstellungen
\usepackage[ngerman]{babel}


\usepackage[
  math-style=ISO,    % ┐
  bold-style=ISO,    % │
  sans-style=italic, % │ ISO-Standard folgen
  nabla=upright,     % │
  partial=upright,   % │
  mathrm=sym,        % ┘
  warnings-off={           % ┐
    mathtools-colon,       % │ unnötige Warnungen ausschalten
    mathtools-overbracket, % │
  },                       % ┘
]{unicode-math}

% traditionelle Fonts für Mathematik
\setmathfont{Latin Modern Math}
% Alternativ zum Beispiel:
%\setmathfont{Libertinus Math}

\setmathfont{XITS Math}[range={scr, bfscr}]
\setmathfont{XITS Math}[range={cal, bfcal}, StylisticSet=1]

% Zahlen und Einheiten
\usepackage[
  locale=DE,                   % deutsche Einstellungen
  separate-uncertainty=true,   % immer Unsicherheit mit \pm
  per-mode=symbol-or-fraction, % / in inline math, fraction in display math
]{siunitx}

% chemische Formeln
\usepackage[
  version=4,
  math-greek=default, % ┐ mit unicode-math zusammenarbeiten
  text-greek=default, % ┘
]{mhchem}

% richtige Anführungszeichen
\usepackage[autostyle]{csquotes}

% schöne Brüche im Text
\usepackage{xfrac}

% Standardplatzierung für Floats einstellen
\usepackage{float}
\floatplacement{figure}{htbp}
\floatplacement{table}{htbp}

% Floats innerhalb einer Section halten
\usepackage[
  section, % Floats innerhalb der Section halten
  below,   % unterhalb der Section aber auf der selben Seite ist ok
]{placeins}

% Seite drehen für breite Tabellen: landscape Umgebung
\usepackage{pdflscape}

% Captions schöner machen.
\usepackage[
  labelfont=bf,        % Tabelle x: Abbildung y: ist jetzt fett
  font=small,          % Schrift etwas kleiner als Dokument
  width=0.9\textwidth, % maximale Breite einer Caption schmaler
]{caption}
% subfigure, subtable, subref
\usepackage{subcaption}

% Grafiken können eingebunden werden
\usepackage{graphicx}

% schöne Tabellen
\usepackage{tabularray}
\UseTblrLibrary{booktabs, siunitx}

% longtblr-Umgebung auf Deutsch
\DefTblrTemplate{contfoot-text}{normal}{Weiter auf der nächsten Seite}
\SetTblrTemplate{contfoot-text}{normal}
\DefTblrTemplate{conthead-text}{normal}{(Fortsetzung)}
\SetTblrTemplate{conthead-text}{normal}

% Verbesserungen am Schriftbild
\usepackage{microtype}

% Literaturverzeichnis
\usepackage[
  backend=biber,
]{biblatex}
% Quellendatenbank
\addbibresource{lit.bib}
\addbibresource{programme.bib}

% Hyperlinks im Dokument
\usepackage[
  german,
  unicode,        % Unicode in PDF-Attributen erlauben
  pdfusetitle,    % Titel, Autoren und Datum als PDF-Attribute
  pdfcreator={},  % ┐ PDF-Attribute säubern
  pdfproducer={}, % ┘
]{hyperref}
% erweiterte Bookmarks im PDF
\usepackage{bookmark}

% Trennung von Wörtern mit Strichen
\usepackage[shortcuts]{extdash}
\usepackage{tikz}
\usepackage[european]{circuitikz}   
% Code Einbindung
\usepackage{xcolor}
\definecolor{LightGray}{gray}{0.9}
\usepackage{minted}
\setminted{
  autogobble,
  breaklines,
  stripnl=true,
  style=tango,
  baselinestretch=1.2,
  bgcolor=LightGray,
  fontsize=\footnotesize,
}

\author{%
  David Lohberg\\%
  \href{mailto:david.lohberg@tu-dortmund.de}{david.lohberg@tu-dortmund.de}%
  \and%
  Felix Michael Bernd Wendeln\\%
  \href{mailto:felix.wendeln@tu-dortmund.de}{felix.wendeln@tu-dortmund.de}%
}
\publishers{TU Dortmund – Fakultät Physik}

\usepackage{expl3}
\usepackage{xparse}

\usepackage{mleftright}
\setlength{\delimitershortfall}{-1sp}

\ExplSyntaxOn

\NewDocumentCommand \I {} { \symup{i} }
\NewDocumentCommand \E {} { \symup{e} }
\NewDocumentCommand \mpi {} { \symup{\pi} }

\NewDocumentCommand \dif {m} { \mathinner{\symup{d} #1} }
\NewDocumentCommand \Dif {m} { \mathinner{\symup{D} #1} }
\NewDocumentCommand \del {m} { \mathinner{\symup{\delta} #1} }
\NewDocumentCommand \Del {m} { \mathinner{\symup{\Delta} #1} }

\let\vaccent=\v % alten Befehl kopieren
\RenewDocumentCommand \v {} % Befehl überschreiben
  {
    \TextOrMath{
      \vaccent % Textmodus
    }{
      \symbf % Mathemodus
    }
  }

\AtBeginDocument{ % wird bei \begin{document} ausgeführt
  % werden sonst wieder von unicode-math überschrieben
  \RenewDocumentCommand \Re {} {\operatorname{Re}}
  \RenewDocumentCommand \Im {} {\operatorname{Im}}
}

\let\ltext=\l
\RenewDocumentCommand \l {}{ \TextOrMath{ \ltext }{ \mleft } }
\let\raccent=\r
\RenewDocumentCommand \r {}{ \TextOrMath{ \raccent }{ \mright } }

\DeclarePairedDelimiter{\abs}{\lvert}{\rvert}
\DeclarePairedDelimiter{\norm}{\lVert}{\rVert}

\DeclarePairedDelimiter{\bra}{\langle}{\rvert}
\DeclarePairedDelimiter{\ket}{\lvert}{\rangle}
% <name> <#arguments> <left> <right> <body>
\DeclarePairedDelimiterX{\braket}[2]{\langle}{\rangle}{#1 \delimsize| #2}


\NewDocumentCommand \OverfullCenter {+m} {
\noindent\makebox[\linewidth]{#1} }

\ExplSyntaxOff


\subject{VERSUCH 350}
\title{Digitale Bauteile}
\date{
  Durchführung: 20.10.2026
  \hspace{3em}
  Abgabe: 03.11.2026
}
\usepackage{subcaption}
\usepackage[section, below]{placeins}
\usepackage[]{caption} % Captions schöner machen
\FloatBarrier
\usepackage{comment}

\begin{document}

\maketitle
\thispagestyle{empty}
\tableofcontents
\newpage
%\input{content/theorie.tex}
%\input{content/durchfuehrung.tex}
%\input{content/auswertung.tex}
%\input{content/diskussion.tex}
\section{Theorie}
\subsection{Logik und Logikpegel}
Wie digitale Bauelemente Signale annehmen, interpretieren und umwandeln, wird durch eine bestimmte Logik und einen Logikpegel festgelegt.
\subsubsection{Logik}
    Die Verknüpfung digitaler Bauelemente und ihrer Signale 
    geschieht durch die logischen Zustände „0“ für „nicht-wahr“ und „1“ für „wahr“.
  \subsubsection{Logikpegel}
    Digitale Schaltungen bestehen aus Transistorschaltungen,
    die je nach verwendeter Halbleitertechnologie auf bestimmte Spannungen reagieren. 
    Beispiele sind dafür TTL oder CMOS, 
    die bei unterschiedlich breiten Voltbereichen ein Low-Level 
    bzw. High-Level Signal ausgeben. 
    Wie diese Signale interpretiert werden,   
    hängt von der Art der verbauten Logik ab. 
    Eine positive Logik übersetzt ein High-Level Signal als 
    logischen Zustand „1“ und ein Low-Level Signal als „0“.
    Bei einer negativen Logik ist dies umgekehrt, 
    also wird ein High-Level Signal zu einer „0“ 
    und ein Low-Level Signal zu einer „1“. \cite{digbau}
\subsection{Gatter-Arten}
  Es gibt verschiedene Arten von Gattern, 
  die Signale in unterschiedlicher Weise umwandeln.\cite{digbau}
  \subsubsection{NOT-Gatter}
    Das NOT-Gatter wirkt als Negation des Signals.
    Es führt bei einem Eingangssignal „0“ zu einem Ausgangssignal „1“ 
    und bei einem Eingangssignal „1“ zu einem Ausgangssignal „0“.
    Wie dieses Gatter in einem Schaltplan dargestellt wird, lässt sich aus der unten stehenden Abbildung erkennen.
    \begin{figure}[H]
       \begin{subfigure}[t]{0.48\textwidth}%
        \centering
        \caption{NOT-Gatter.\cite{tikz}}
        \vspace{15pt}
        \begin{tikzpicture}[transform shape]
	        \node[european not port] at (2, 0.4){};
	        \node[notcirc] at (2.62, 0.4){};
	        \node[notcirc] at (-0.15, 0.4){};
          \node at (-0.6,0.4){$E$};
          \node at (3, 0.45){$A$};
          \draw (-0.075, 0.4) -- (0.23, 0.4);
	        \draw (2.195, 0.4) -- (2.545, 0.4);
        \end{tikzpicture}
      \end{subfigure}
      \hfill
      \begin{subfigure}[t]{0.48\textwidth}
        \centering
        \vspace{0pt}
        \caption{Wahrheitstabelle NOT-Gatter.\cite{digbau}}
        \begin{tabular}{c c}
          \toprule
          $E$ & $A$\\
          \midrule
          0 & 1 \\
          1 & 0 \\
          \bottomrule
        \end{tabular}
      \end{subfigure}
    \end{figure}

 \subsubsection{OR-Gatter}
    Das OR-Gatter vergleicht mehrere Eingangssignale miteinander.
    Ist mindestens ein Eingangssignal eine „1“ wird auch das Ausgangssignal eine „1“.
    Nur wenn alle Eingangssignale „0“ sind, wird auch eine „0“ ausgegeben.
     Das OR-Gatter wird wie in der unteren Abbildung in Schaltplänen dargestellt.
         \begin{figure}[H]
       \begin{subfigure}[t]{0.48\textwidth}%
        \centering
        \caption{OR-Gatter.\cite{tikz}}
        \vspace{15pt}
        \begin{tikzpicture}[transform shape]
	        \node[european or port] at (2, 0.4){};
	        \node[notcirc] at (2.62, 0.4){};
	        \node[notcirc] at (-0.15, 0.63){};
          \node[notcirc] at (-0.15, 0.17){};
          \node at (-0.6,0.63){$E_1$};
          \node at (-0.6, 0.17){$E_2$};
          \node at (3, 0.45){$A$};
          \draw (-0.075, 0.63) -- (0.25, 0.63);
           \draw (-0.075, 0.17) -- (0.25, 0.17);
	        \draw (2.195, 0.4) -- (2.545, 0.4);
        \end{tikzpicture}
      \end{subfigure}
      \hfill
      \begin{subfigure}[t]{0.48\textwidth}
        \centering
        \vspace{0pt}
        \caption{Wahrheitstabelle OR-Gatter.\cite{digbau}}
        \begin{tabular}{c c c}
          \toprule
          $E_1$ & $E_2$ & $A$\\
          \midrule
          0 & 0 & 1\\
          0 & 1  & 0\\
          1 & 0  & 0\\
          1 & 1 & 0 \\
          \bottomrule
        \end{tabular}
      \end{subfigure}
    \end{figure}
  \subsubsection{AND-Gatter}
    Das AND-Gatter vergleicht mehrere Eingangssignale miteinander 
    und gibt dementsprechend eine „1“ oder eine „0“ aus. 
    Nur wenn alle Eingangssignale eine „1“ sind gibt das Gatter das Ausgangssignal „1“.
    Wenn mindestens ein Eingangssignal „0“ ist, wird stattdessen das Signal „0“ ausgegeben.
    Die Darstellung innerhalb eines Schaltplans ist in einer Darstellung unten gezeigt.
      \begin{figure}[H]
       \begin{subfigure}[t]{0.48\textwidth}%
        \centering
        \caption{AND-Gatter.\cite{tikz}}
        \vspace{15pt}
        \begin{tikzpicture}[transform shape]
	        \node[european and port] at (2, 0.4){};
	        \node[notcirc] at (2.62, 0.4){};
	        \node[notcirc] at (-0.15, 0.63){};
          \node[notcirc] at (-0.15, 0.17){};
          \node at (-0.6,0.63){$E_1$};
          \node at (-0.6, 0.17){$E_2$};
          \node at (3, 0.45){$A$};
          \draw (-0.075, 0.63) -- (0.25, 0.63);
           \draw (-0.075, 0.17) -- (0.25, 0.17);
	        \draw (2.195, 0.4) -- (2.545, 0.4);
        \end{tikzpicture}
      \end{subfigure}
      \hfill
      \begin{subfigure}[t]{0.48\textwidth}
        \centering
        \vspace{0pt}
        \caption{Wahrheitstabelle AND-Gatter.\cite{digbau}}
        \begin{tabular}{c c c}
          \toprule
          $E_1$ & $E_2$ & $A$\\
          \midrule
          0 & 0 & 0\\
          0 & 1  & 0\\
          1 & 0  & 0\\
          1 & 1 & 1 \\
          \bottomrule
        \end{tabular}
      \end{subfigure}
    \end{figure}
  \subsubsection{NAND-Gatter}
    Das NAND-Gatter kombiniert die Wirkungsweise eines NOT-Gatters mit einem AND-Gatter. 
    Sind alle Eingangssignale „1“ wird das Ausgangssignal eine „0“. Ist dies nicht der Fall,
     wird eine „1“ ausgegeben.
    Die Darstellung dieses Gatters ist auch in Schaltplänen eine Kombination eines AND- und NOT-Gatters, wie es sich aus der unten stehenden Abbildung entnehmen lässt.
              \begin{figure}[H]
       \begin{subfigure}[t]{0.48\textwidth}%
        \centering
        \caption{Or-Gatter.\cite{tikz}}
        \vspace{15pt}
        \begin{tikzpicture}[transform shape]
	        \node[european nand port] at (2, 0.4){};
	        \node[notcirc] at (2.62, 0.4){};
	        \node[notcirc] at (-0.15, 0.63){};
          \node[notcirc] at (-0.15, 0.17){};
          \node at (-0.6,0.63){$E_1$};
          \node at (-0.6, 0.17){$E_2$};
          \node at (3, 0.45){$A$};
          \draw (-0.075, 0.63) -- (0.25, 0.63);
           \draw (-0.075, 0.17) -- (0.25, 0.17);
	        \draw (2.195, 0.4) -- (2.545, 0.4);
        \end{tikzpicture}
      \end{subfigure}
      \hfill
      \begin{subfigure}[t]{0.48\textwidth}
        \centering
        \vspace{0pt}
        \caption{Wahrheitstabelle OR-Gatter.\cite{digbau}}
        \begin{tabular}{c c c}
          \toprule
          $E_1$ & $E_2$ & $A$\\
          \midrule
          0 & 0 & 0\\
          0 & 1  & 1\\
          1 & 0  & 1\\
          1 & 1 & 1 \\
          \bottomrule
        \end{tabular}
      \end{subfigure}
    \end{figure}
\subsection{Anwendungen}
  \subsubsection{NAND-Gatter-Komplexe}
    Durch Verknüpfung von mehreren NAND-Gattern lassen sich andere Arten von Gattern bauen. 
    Wird an die Eingänge eines NAND-Gatters das selbe Signal gesendet, 
    wirkt das NAND-Gatter wie ein NOT-Gatter.Wie diese Anwendung eines NAND-Gatters zu schalten ist und logisch funktioniert, ist in der unteren Abbildung dargestellt und über Boolsche Algebra gezeigt.
       \begin{figure}[H]
       \begin{subfigure}[t]{0.48\textwidth}%
        \centering
        \caption{NOT-Gatter aus NAND-Gatter.\cite{tikz}}
        \vspace{15pt}
        \begin{tikzpicture}[transform shape]
	        \node[european nand port] at (2, 0.4){};
	        \node[notcirc] at (2.62, 0.4){};
	        \node[notcirc] at (-0.15, 0.4){};
          \node at (-0.6,0.4){$A$};
          \node at (3, 0.45){$\overline{A}$};
          \draw (-0.075, 0.4) -- (0.23, 0.4);
	        \draw (0.23, 0.635) -- (0.23, 0.175);
	        \draw (2.195, 0.4) -- (2.545, 0.4);
        \end{tikzpicture}
      \end{subfigure}
      \hfill
      \begin{subfigure}[t]{0.48\textwidth}
        \centering
        \vspace{0pt}
        \caption{Boolsche Algebra.}
        \begin{equation}
          \overline{A} = \overline{A \wedge B}
        \end{equation}
      \end{subfigure}
    \end{figure}
    \noindent
    Schließt man ein NAND-Gatter, als NOT-Gatter, 
    an ein anderes NAND-Gatter wird das erstere negiert und wirkt somit wie ein AND-Gatter. 
    Die verwendete Kombination an Gattern wird im folgenden Schaltbild abgebildet und die Logik erneut über Boolsche Algebra erklärt.
    \begin{figure}
      \begin{subfigure}[t]{0.48\textwidth}%
        \centering
        \caption{UND-Gatter aus NAND-Gatter.\cite{tikz}}
        \vspace{15pt}
        \begin{tikzpicture}[transform shape]
	        \node[european nand port] at (2, 0.75){};
	        \node[european nand port] at (4.304, 0.75){};
	        \node[notcirc] at (0.15, 0.977){};
	        \node[notcirc] at (0.15, 0.522){};
	        \node[notcirc] at (5, 0.75){};
          \node at (-0.3,0.977){$A$};
          \node at (-0.3,0.522){$B$};
          \node at (5, 1.15){$A \wedge B$};
	        \draw (2.196, 0.75) -| (2.54, 0.522);
	        \draw (2.54, 0.75) -- (2.54, 0.977);
	        \draw (4.5, 0.75) -- (4.914, 0.75);
        \end{tikzpicture}
      \end{subfigure}
      \hfill
      \begin{subfigure}[t]{0.48\textwidth}
        \centering
        \vspace{0pt}
        \caption{Boolsche Algebra.}
        \begin{equation}
          A \wedge B = \overline{ (\overline{(A \wedge B)}) }
        \end{equation}
      \end{subfigure}
    \end{figure}

    %Verbindet man zwei, wie NICHT-Gatter wirkende,
    %NAND-Gatter an ein drittes NAND-Gatter, wirken diese wie ein OR-Gatter.
    %\begin{figure}[H]
    %  \begin{subfigure}[t]{0.48\textwidth}%
    %    \centering
    %    \caption{OR-Gatter aus NAND-Gatter\cite{tikz}}
    %    \vspace{5pt}
    %    \begin{tikzpicture}[transform shape]
	        %\node[european nand port] at (2, 0.4){};
	        %\node[european nand port] at (4.3, 1){};
          %\node[european nand port] at (2,1.6){};
	        %\node[notcirc] at (-0.15, 1.6){};
	        %\node[notcirc] at (-0.15, 0.4){};
	        %\node[notcirc] at (5, 1){};
          %\node at (-0.6,1.6){$A$};
          %\node at (-0.6,0.4){$B$};
          %\node at (5, 1.35){$A \vee B$};
	        %\draw (-0.075, 1.6) -- (0.23, 1.6);
	        %\draw (0.23, 1.835) -- (0.23, 1.365);
          %\draw (-0.075, 0.4) -- (0.23, 0.4);
	        %\draw (0.23, 0.635) -- (0.23, 0.175);
	        %\draw (2.195, 0.4) -- (2.545, 0.4);
	        %\draw (2.545, 0.4) -- (2.545, 0.765);
          %\draw (2.195, 1.6) -- (2.545, 1.6);
	        %\draw (2.545, 1.6) -- (2.545, 1.235);
	        %\draw (4.5, 1) -- (4.915, 1);
        %\end{tikzpicture}
      %\end{subfigure}
      %\hfill
      %\begin{subfigure}[t]{0.48\textwidth}
       % \centering
        %\vspace{0pt}
        %\caption{Boolsche Algebra}
        %\begin{equation}
        %  A \vee B = \overline{(\overline{A} \wedge \overline{B})}
        %\end{equation}
      %\end{subfigure}
    %\end{figure}


    %Auch lässt sich aus drei NAND-Gattern mit zwei Eingängen eine Schaltung die wie ein NAND-Gatter mit drei Eingängen wirkt bauen. 
    %Dafür schaltet man zwei NAND-Gatter, die wie ein AND-Gatter wirken, an ein drittes NAND-Gatter, 
    %welches an einen dritten Eingang angeschlossen ist.

    %\begin{figure}
    %  \begin{subfigure}[t]{0.48\textwidth}%
    %    \centering
    %    \caption{3-NAND aus NAND-Gatter\cite{tikz}}
    %    \vspace{5pt}
    %    \begin{tikzpicture}[transform shape]
	        % Paths, nodes and wires:
	        %\node[european nand port] at (7, 0.627){};
	        %\node[european nand port] at (2, 2.2){};
          %\node[european nand port] at (4.5,2.2){};
	        %\node[notcirc] at (0.15, 2.427){};
	        %\node[notcirc] at (0.15, 0.4){};
          %\node[notcirc] at (0.15, 1.973){};
	        %\node[notcirc] at (8, 0.627){};
          %\node at (-0.3,1.93){$B$};
          %\node at (-0.3,0.4){$C$};
          %\node at (-0.3, 2.47){$A$};
          %\node at (8, 1){$\overline{A \wedge B \wedge C}$};
          %\draw (7, 0.627) -- (7.9, 0.627);
	        %\draw (0.25 , 0.4) -- (5.5, 0.4);
          %\draw (2, 2.2) -- (2.75, 2.2);
	        %\draw (2.74, 1.973) -- (2.74, 2.427);
	        %\draw (4.69, 2.2) -- (4.69, 0.854);
	        %\draw (4.69, 0.854) -- (5.5, 0.854);
        %\end{tikzpicture}
      %\end{subfigure}
      %\hfill
      %\begin{subfigure}[t]{0.48\textwidth}
        %\centering
        %\vspace{0pt}
        %\caption{Boolsche Algebra}
        %\begin{equation}
          %\overline{A \wedge B \wedge C} = \overline{\overline{\overline{(A \wedge B)}}\wedge C}
        %\end{equation}
      %\end{subfigure}
    %\end{figure}
  \subsubsection{RS-Flip-Flop}
     Ein RS-Flip-Flop ist eine Schaltung, die auf Zustandsänderungen bei dem Eingang reagiert,
     also eine zeitabhängige Schaltung. 
     Sie kann aus vier NAND-Gattern gebaut werden 
     (Mindestens zwei sind für die Logik eines Flip-Flops nötig). 
     Die Schaltung hat zwei Eingänge, die als „S“ bzw. „R“ für „Setzen“ bzw. „Rücksetzen“ bezeichnet werden. 
     Weiterhin hat sie zwei gegensätzliche Ausgänge „Q“ und „$\overline{Q}$“. Das Schaltbild dieses Flip-Flops ist in der unten stehenden Abbildung zu erkennen.

     \noindent
     Das Schaltbild dieses Flip-Flops ist in der unten stehenden Abbildung zu erkennen.
     Wenn „Q“ auf „1“ gesetzt ist, ist „$\overline{Q}$“ auf „0“ gesetzt und umgekehrt. 
     Das wechseln dieser Zustände wird durch das setzen einer „1“ auf dem „S“ bzw. „R“ Eingang. 
     Bei einer „1“ bei dem Eingang „S“ und Ausgang „$\overline{Q}$“ wird letzterer auf „0“ und „Q“ auf „1“ gesetzt. 
     Bei einer „1“ bei dem Eingang „R“ und Ausgang „Q“ wird dieser auf „0“ und „$\overline{Q}$“ auf „1“ gesetzt. 
     Ohne neues Signal (0-0) wird das vorherige Signal erhalten. 
     Werden beide Eingänge gleichzeitig auf „1“ gesetzt, kommt es zu undefiniertem Verhalten. 
     Es ist unvorhersehbar durch beispielsweise variierende Halbleiterstrukturen der Gatter. 
     Die Aus- und Eingabe von Daten ist durch Taktsignale steuerbar.
    Das Verhalten dieses Flip-Flops bei unterschiedlichen Signalen wird in der folgenden Wahrheitstabelle verdeutlicht..\cite{digbau}
    \begin{figure}[H]
      \begin{subfigure}[t]{0.48\textwidth}%
        \centering
        \caption{RS-Flip-Flop.\cite{tikz}}
        \vspace{0pt}
        \begin{tikzpicture}[transform shape, scale = 1]
    	    \node[european nand port] at (2.5, 5){};
    	    \node[european nand port] at (5.264, 4.978){};
	        \node[european nand port] at (5.264, 2.727){};
	        \draw (0.486, 4.977) -- (-0.014, 4.977) -- (0.486, 4.977) -| (0.736, 5.227);
	        \draw (0.736, 4.977) -- (0.736, 4.772);
	        \node[jump crossing] at (3.5, 4){};
	        \node[shape=rectangle, minimum width=-0.035cm, minimum height=-0.035cm] at (0, -0){};
          \draw (5.46, 2.727) -- (5.46, 3.5);
	        \draw (3.5, 4.75) -| (3.5, 4.14);
	        \draw (5.46, 3.5) |- (3.5, 3.5) -| (3.5, 3.86);
	        \draw (3.5, 2.955) -| (2.75, 3) -| (2.75, 4) -- (3.36, 4);
	        \draw (3.64, 4) -| (5.46, 4.978);
	        \draw (5.46, 4.978) -| (7, 5);
	        \draw (5.46, 2.727) -- (7, 2.727);
	        \node[ocirc] at (7, 4.978){};
	        \node[ocirc] at (7, 2.727){};
	        \node[ocirc] at (-0.014, 4.978){};
	        \draw (2.696, 5) -| (3.5, 5.205);
	        \draw (3.5, 2.5) -- (2.75, 2.5);
	        \node[european nand port] at (2.554, 2.5){};
	        \draw (0.54, 2.477) -- (0.04, 2.477) -- (0.54, 2.477) -| (0.79, 2.727);
	        \draw (0.79, 2.477) -- (0.79, 2.272);
	        \node[ocirc] at (0.04, 2.477){};
          \node at (0,2.8) {R};
          \node at (0,5.3) {S};
          \node at (7,5.3) {Q};
          \node at (7,3.2) {$\overline{Q}$};
        \end{tikzpicture}
      \end{subfigure}
      \hfill
      \begin{subfigure}[t]{0.48\textwidth}
        \centering
        \vspace{0pt}
        \caption{Wahrheitstabelle vom Flip-Flop.\cite{digbau}}
        \begin{tabular}{c c c c c}
          \toprule
          $S$ & $R$ & $Q$ & $\overline{Q}$ & Bedeutung\\
          \midrule
          1 & 0 & 1 & 0 & Setzen\\
          0 & 1 & 0 & 1 & Rücksetzen\\
          0 & 0 & 1 & 0 & Speichern\\
          1 & 1 & 1 & 1 & Undefiniert\\
          \bottomrule
        \end{tabular}
      \end{subfigure}
    \end{figure}

  \subsubsection{D-Flip-Flop}
    Ein D-Flip-Flop ist ein taktgesteuerter RS-Flip-Flop mit einem zusätzlichem NOT-Gatter. 
    Dieses verhindert durch setzen von „S“ = „$\overline{R}$“ undefiniertes Verhalten.
    Wie das Schaltbild eines D-Flip-Flops aussieht und wie es sich logisch verhält, lässt sich aus der unteren Abbildung und der Wahrheitstabelle ziehen.
    \begin{figure}[H]
      \begin{subfigure}[t]{0.56\textwidth}%
        \centering
        \caption{D-Flip-Flop.\cite{tikz}}
        \vspace{0pt}
        \begin{tikzpicture}[transform shape, scale = 1]
    	    \node[european nand port] at (2.554, 5){};
    	    \node[european nand port] at (5.264, 4.978){};
	        \node[european nand port] at (5.264, 2.727){};
	        \node[jump crossing] at (3.5, 4){};
	        \node[shape=rectangle, minimum width=-0.035cm, minimum height=-0.035cm] at (0, -0){};
          \draw (5.46, 2.727) -- (5.46, 3.5);
	        \draw (3.5, 4.75) -| (3.5, 4.14);
	        \draw (5.46, 3.5) |- (3.5, 3.5) -| (3.5, 3.86);
	        \draw (3.5, 2.955) -| (2.75, 3) -| (2.75, 4) -- (3.36, 4);
	        \draw (3.64, 4) -| (5.46, 4.978);
	        \node[ocirc] at (-1.8, 5.227){};
	        \draw (2.696, 5) -| (3.5, 5.205);
          \node [jump crossing] at (-0.96, 3.75){};
	        \draw (3.5, 2.5) -- (2.75, 2.5);
	        \node[european nand port] at (2.554, 2.5){};
          \node[european not port] at (0.8, 2.272){};
	        \draw (-1.735, 5.227) -- (0.85, 5.227);
          \draw (-0.96,5.227) -- (-0.96,2.272);
          \draw (-1.735, 3.75) -- (-1.03, 3.75);
          \draw (0.79, 2.73) -- (0.79, 4.773);
          \draw (-0.91, 3.75) -- (0.79, 3.75);
	        \node[ocirc] at (-1.8, 3.75){};
          \node at (-1.8,4) {C};
          \node at (-1.8,5.477) {D};
          \node [ocirc] at (6, 2.727){};
          \node [ocirc] at (6, 4.978) {};
          \draw (5.46, 2.727) -- (5.95, 2.727);
          \draw (5.46, 4.978) -- (5.95, 4.978);
          \node at (6,5.3) {Q};
          \node at (6,3.2) {$\overline{Q}$};
        \end{tikzpicture}
      \end{subfigure}
      \hfill
      \begin{subfigure}[t]{0.4\textwidth}
        \centering
        \vspace{0pt}
        \caption{Wahrheitstabelle vom D-Flip-Flop.\cite{digbau}}
        \begin{tabular}{c c c c c}
          \toprule
          $C$ & $D$ & $Q$ & $\overline{Q}$ & Bedeutung\\
          \midrule
          1 & 1 & 1 & 0 & Setzen\\
          1 & 0 & 0 & 1 & Rücksetzen\\
          0 & 1 & 1 & 0 & Speichern\\
          0 & 0 & 1 & 0 & Speichern\\
          \bottomrule
        \end{tabular}
      \end{subfigure}
    \end{figure}
    \noindent
    Bei dem Anlegen einer periodischen Rechteckspannung am Eingang „C“ wird die Ausgangsfrequenz zu der halben Eingangsfrequenz. 
    Der Flip-Flop wirkt als Frequenzteiler.\cite{digbau}
    Es lassen sich auch mehrere solcher Frequenzteiler hintereinander schalten. Jedes dieser Flip-Flops halbiert die Frequenz, wodurch sich eine maximale Frequenzteilung von $2^n$, bei $n$ Flip-Flops, erreichen lässt. In dem folgenden Impulsdiagramm wird der Unterschied der Frequenz vor und nach Erreichen des Flip-Flops verdeutlicht.
    %Plot Frequenz teiler
    \begin{figure}
\centering
\includegraphics[width=\textwidth]{frequenzteiler.pdf}
\caption{Frequnezteiler\cite{digbau,matplotlib}}
\end{figure}
    
    \begin{comment}
    \subsection{Weitere Bauteile}
  \subsubsection{Binärzähler}
    Ein Binärzähler besteht aus mehreren D-Flip-Flops, 
    die als Frequenzteiler nacheinandergeschaltet werden. 
    Diese werden mit $2^n$ nummeriert. Mit Jumpern lässt sich die gewollte Frequenzteilung bestimmen. 
    So lassen sich Binärzahlen zählen. Ein High-Pegel an dem Reset Eingang führt zum Zurücksetzen des Zählers. 
    Durch zusätzliche Gatter kann der Zähler direkt beim erreichen bestimmter Zahlen zurückgesetzt werden.\cite{digbau}
  \subsubsection{Quarzoszillator}
    Ein Quarzoszillator erzeugt eine Frequenz, die mit Frequenzteilern auf die gewollte Frequenz heruntergebrochen wird.\cite{digbau}
  \subsubsection{Dekadenzähler}
    Ein Dekadenzähler ist ein Zähler mit einem Takteingang und mehreren Ausgängen, die mit LEDs verbunden sind. 
    Es wird nur einer der Ausgänge gesetzt.\cite{digbau} 
    Durch das Verbinden eines Ausgangs mit dem Reset Eingang werden nur die LEDs bis zu diesem Ausgang verwendet. 
    Ein High-Pegel verbunden mit dem Reset Eingang führt auch hier zum Zurücksetzen des Zählers.\cite{digbau}
  \subsubsection{7-Segment-Anzeige}
    Eine 7-Segment-Anzeige setzt Binärzahlen über 7-Segmentcode um und
    zeigt sie als Dezimalzahl durch Aufleuchten der jeweiligen Segmente an.\cite{digbau}
  \subsubsection{Piezo}
    Ein Piezo erzeugt Signaltöne in einem breiten Frequenz- und Lautstärkebereich. \cite{digbau}
\end{comment}

\printbibliography{}

%\input{content/minted.tex}


\end{document}
